\documentclass[11pt]{article}
\usepackage[utf8]{inputenc}
\usepackage{amsmath,amssymb,booktabs,geometry,verbatim}
\usepackage{CJKutf8}
\geometry{margin=1in}
\usepackage[hidelinks]{hyperref}

\title{CSP-S 2024 (Senior)}
\date{2024-10-26, 14:30--18:30 (4 hours)}

\begin{document}
\maketitle

\noindent
Translated from the official Chinese statements of the second round of the
2024 CCF Certified Software Professional (CSP-J/S 2024), Senior level.
The problems are given in contest order.
File names in parentheses are the official problem code names.
Each problem reads from \texttt{<name>.in} and writes to \texttt{<name>.out}.
Merged cells of the original subtask tables are reproduced by an entry
in the first row of the merged block and empty cells below it.

\clearpage
\section*{Problem 1: Duel (\begin{CJK}{UTF8}{gbsn}决斗\end{CJK})
  \hfill{\normalsize\href{https://qoj.ac/problem/10194}{QOJ 10194}}}

\noindent
Problem code name: \texttt{duel}.
Time limit: 1 s.
Memory limit: 512 MB.
Total points: 100.

\subsection*{Problem description}

Today is little Q's birthday, and he received $n$ cards as a present.
These cards belong to the popular game ``Duel Monsters''.
The $i$-th card represents a monster whose attack power is $r_i$ and whose
defense power is also $r_i$.

A game consists of several rounds.
In each round, little Q chooses a monster $i$ and \emph{another} monster
$j$ ($i \neq j$), and lets monster $i$ attack monster $j$.
If the attack power of monster $i$ is less than or equal to the defense
power of monster $j$, nothing happens.
Otherwise, the defense of monster $j$ is broken, and monster $j$ leaves the
game and no longer takes part in the rest of it.
A monster can launch an attack \emph{at most once} in the whole game.
The game ends when every monster that has not left the game has already
launched an attack.

Little Q wants to decide an order of attacks such that, when the game
ends, the number of monsters that have not left the game is as small as
possible.

\subsection*{Input}

Read the input from the file \texttt{duel.in}.

The first line of the input contains a positive integer $n$, the number of
cards.

The second line of the input contains $n$ positive integers, the $i$-th of
which is the attack and defense power $r_i$ of the $i$-th monster.

\subsection*{Output}

Write the output to the file \texttt{duel.out}.

Output one line containing one integer, the minimum possible number of
monsters that have not left the game when the game ends.

\subsection*{Sample 1 input}

\begin{verbatim}
5
1 2 3 1 2
\end{verbatim}

\subsection*{Sample 1 output}

\begin{verbatim}
2
\end{verbatim}

\subsection*{Explanation of sample 1}

One optimal plan is the following.
In the first round, let monster 2 attack monster 1.
In the second round, let monster 5 attack monster 4.
In the third round, let monster 3 attack monster 5.
Now every monster that has not left the game has attacked, so the game
ends.
It can be shown that no attack order does better.

\subsection*{Sample 2 input}

\begin{verbatim}
10
136 136 136 2417 136 136 2417 136 136 136
\end{verbatim}

\subsection*{Sample 2 output}

\begin{verbatim}
8
\end{verbatim}

\subsection*{Sample 3}

See the files \texttt{duel/duel3.in} and \texttt{duel/duel3.ans} in the
contestant's directory.
On QOJ these are the downloadable attachment files \texttt{ex\_duel3.in}
and \texttt{ex\_duel3.ans}.

This sample satisfies $\forall 1 \le i \le n,\ r_i \le 2$.

\subsection*{Sample 4}

See the files \texttt{duel/duel4.in} and \texttt{duel/duel4.ans} in the
contestant's directory.
On QOJ these are the downloadable attachment files \texttt{ex\_duel4.in}
and \texttt{ex\_duel4.ans}.

\subsection*{Constraints}

For all test data it is guaranteed that $1 \le n \le 10^5$ and
$1 \le r_i \le 10^5$.

\begin{center}
\begin{tabular}{|c|c|c|c|}
\hline
Test cases & $n$ & $r_i$ & Special property \\
\hline
$1 \sim 4$ & $\le 10$ & $\le 10^5$ & none \\
\cline{1-3}
$5 \sim 10$ & $\le 10^5$ & $\le 2$ & \\
\hline
$11 \sim 15$ & $\le 30$ & $\le 10^5$ & special property A \\
\cline{1-2}\cline{4-4}
$16 \sim 20$ & $\le 10^5$ & & none \\
\hline
\end{tabular}
\end{center}

Special property A: it is guaranteed that every $r_i$ is generated
independently and uniformly at random from its possible range of values.

\clearpage
\section*{Problem 2: Speeding Detection
(\begin{CJK}{UTF8}{gbsn}超速检测\end{CJK})
  \hfill{\normalsize\href{https://qoj.ac/problem/10195}{QOJ 10195}}}

\noindent
Problem code name: \texttt{detect}.
Time limit: 2 s.
Memory limit: 512 MB.
Total points: 100.

\subsection*{Problem description}

Little D has just joined the traffic administration of a certain country.
His first task is to be in charge of speeding detection on a north--south
main road of the country whose length is $L$.
To test little D, his supervisor first asks him to solve a simplified
scenario.

This weekend, $n$ cars are expected to appear on the main road.
The $i$-th car enters the main road at the position at distance $d_i$ from
its southern end, and drives northwards with uniform acceleration, with
initial speed $v_i$ and acceleration $a_i$.
We only consider cars driving from south to north, so $v_i > 0$, but $a_i$
may be positive, negative or zero.
When a car reaches the northern end of the main road (that is, the position
at distance $L$ from the southern end) or its speed drops to $0$ (which can
only happen when $a_i < 0$), we consider that the car leaves the main road.

There are $m$ speed cameras on the main road.
The $j$-th speed camera is at the position at distance $p_j$ from the
southern end of the main road.
Every speed camera can be set to on or off.
When a car passes a speed camera that is on, if the instantaneous speed of
the car \emph{exceeds} the speed limit $V$ of the road, the car is judged
to be speeding.
Note that when a car enters or leaves the main road, if there is a speed
camera that is on at the corresponding position, that speed camera also
measures the speed of the car.

The supervisor first wants to know how many of these $n$ cars would be
judged to be speeding if all speed cameras were on.

Second, to save energy, the department wants to switch off some of the
speed cameras.
However, they do not want to miss any speeding car.
That is, if one of the $n$ cars is judged to be speeding when all speed
cameras are on, they want it to still be judged to be speeding after some
speed cameras have been switched off.
The supervisor also wants to know the maximum number of speed cameras that
can be switched off under this condition.

Since $n$ is very large, the supervisor allows little D to solve these two
problems by programming, so little D turns to you.

If you are not familiar with acceleration, little D has kindly provided the
formulas about acceleration in the ``Hint'' section of this problem.

\subsection*{Input}

Read the input from the file \texttt{detect.in}.

\emph{This problem has multiple test cases.}

The first line of the input contains a positive integer $T$, the number of
test cases.

Then $T$ test cases follow, each in the following format.

The first line contains four integers $n, m, L, V$: the number of cars,
the number of speed cameras, the length of the main road and the speed
limit of the road, respectively.

Then $n$ lines follow.
The $i$-th of them contains three integers $d_i, v_i, a_i$ describing a
car.

The last line contains $m$ integers $p_1, p_2, \cdots, p_m$ describing the
positions of all speed cameras on the road.

\subsection*{Output}

Write the output to the file \texttt{detect.out}.

For each test case, output one line containing two integers.
The first integer is the number of cars judged to be speeding when all
speed cameras are on.
The second integer is the maximum number of speed cameras that can be
switched off without missing any speeding car.

\subsection*{Sample 1 input}

\begin{verbatim}
1
5 5 15 3
0 3 0
12 4 0
1 1 4
5 5 -2
6 4 -4
2 5 8 9 15
\end{verbatim}

\subsection*{Sample 1 output}

\begin{verbatim}
3 3
\end{verbatim}

\subsection*{Explanation of sample 1}

In this test case the main road has length $15$ and the speed limit is $3$.
There is one speed camera at each of the positions at distance
$2, 5, 8, 9, 15$ from the southern end.
\begin{itemize}
\item The first car enters at the southern end and drives at the constant
speed $3$.
This car is never speeding on the whole road.
\item The second car enters at the position at distance $12$ from the
southern end and drives at the constant speed $4$.
When it leaves the main road at the northern end, it is judged to be
speeding by the speed camera at distance $15$ from the southern end.
\item The third car enters at the position at distance $1$ from the
southern end, with initial speed $1$ and acceleration $4$.
After it has travelled a distance of $\frac{3^2 - 1^2}{2 \times 4} = 1$,
that is, when it reaches position $2$, its speed becomes $3$, and it is
speeding from then on.
Therefore this car is judged to be speeding by every speed camera except
the one at distance $2$ from the southern end.
\item The fourth car enters at the position at distance $5$ from the
southern end, with initial speed $5$ and acceleration $-2$.
After it has travelled a distance of
$\frac{3^2 - 5^2}{2 \times (-2)} = 4$, that is, when it reaches position
$9$, its speed becomes $3$.
Therefore this car is speeding while its distance from the southern end is
in $[5, 9)$, and it is judged to be speeding by the two speed cameras at
distance $5$ and $8$ from the southern end.
\item The fifth car enters at the position at distance $6$ from the
southern end, with initial speed $4$ and acceleration $-4$.
After it has travelled a distance of
$\frac{3^2 - 4^2}{2 \times (-4)} = \frac{7}{8}$, that is, when it reaches
position $6\frac{7}{8}$, its speed becomes $3$.
Therefore this car is speeding while its distance from the southern end is
in $[6, 6\frac{7}{8})$, but there is no speed camera in this interval, so
it is not judged to be speeding.
\end{itemize}
Hence the second, third and fourth cars are judged to be speeding, and the
first number of the output is $3$.

We can switch off the three speed cameras at distance $2, 8, 9$ from the
southern end and keep the two at $5$ and $15$.
Then the three cars that were judged to be speeding before are still judged
to be speeding.
It can be shown that there is no better plan, so the second number of the
output is $3$.

\subsection*{Sample 2}

See the files \texttt{detect/detect2.in} and \texttt{detect/detect2.ans}
in the contestant's directory.
On QOJ these are the downloadable attachment files
\texttt{ex\_detect2.in} and \texttt{ex\_detect2.ans}.

This sample satisfies $n, m \le 10$.

\subsection*{Sample 3}

See the files \texttt{detect/detect3.in} and \texttt{detect/detect3.ans}
in the contestant's directory.
On QOJ these are the downloadable attachment files
\texttt{ex\_detect3.in} and \texttt{ex\_detect3.ans}.

This sample satisfies special property A, and its first ten test cases
satisfy $n, m \le 3000$.

\subsection*{Sample 4}

See the files \texttt{detect/detect4.in} and \texttt{detect/detect4.ans}
in the contestant's directory.
On QOJ these are the downloadable attachment files
\texttt{ex\_detect4.in} and \texttt{ex\_detect4.ans}.

This sample satisfies special property B, and its first ten test cases
satisfy $n, m \le 3000$.

\subsection*{Sample 5}

See the files \texttt{detect/detect5.in} and \texttt{detect/detect5.ans}
in the contestant's directory.
On QOJ these are the downloadable attachment files
\texttt{ex\_detect5.in} and \texttt{ex\_detect5.ans}.

This sample satisfies special property C, and its first ten test cases
satisfy $n, m \le 3000$.

\subsection*{Constraints}

For all test data it is guaranteed that:
\begin{itemize}
\item $1 \le T \le 20$;
\item $1 \le n, m \le 10^5$, $1 \le L \le 10^6$, $1 \le V \le 10^3$;
\item $0 \le d_i < L$, $1 \le v_i \le 10^3$, $|a_i| \le 10^3$;
\item $0 \le p_1 < p_2 < \cdots < p_m \le L$.
\end{itemize}

\begin{center}
\begin{tabular}{|c|c|c|}
\hline
Test case & $n, m \le$ & Special property \\
\hline
1 & $10$ & none \\
\cline{1-2}
2 & $20$ & \\
\hline
3 & $3{,}000$ & A \\
\cline{1-2}
4 & $10^5$ & \\
\hline
5 & $3{,}000$ & B \\
\cline{1-2}
6 & $10^5$ & \\
\hline
7 & $3{,}000$ & C \\
\cline{1-2}
8 & $10^5$ & \\
\hline
9 & $3{,}000$ & none \\
\cline{1-2}
10 & $10^5$ & \\
\hline
\end{tabular}
\end{center}

Special property A: it is guaranteed that $a_i = 0$.

Special property B: it is guaranteed that $a_i > 0$.

Special property C: it is guaranteed that $a_i < 0$, and that no car
leaves the main road at its northern end.

\subsection*{Hint}

The definitions and formulas related to acceleration are as follows.
\begin{itemize}
\item \emph{Uniformly accelerated motion} is motion in which the
acceleration of the object stays constant during the motion, that is, the
change of speed per unit of time is constant.
\item If a car has initial speed $v_0$ and acceleration $a \neq 0$ and
moves with uniform acceleration, then after time $t$ its speed is
$v_1 = v_0 + a \times t$, and its displacement (that is, the distance
travelled) is $s = v_0 \times t + 0.5 \times a \times t^2$.
\item If a car has initial speed $v_0$ and acceleration $a \neq 0$ and
moves with uniform acceleration, then when its displacement (that is, the
distance travelled) is $s$, its instantaneous speed is
$\sqrt{v_0^2 + 2 \times a \times s}$.
\item If a car has initial speed $v_0$ and acceleration $a \neq 0$, then
when its displacement (that is, the distance travelled) is
$\frac{v_1^2 - v_0^2}{2a}$, its instantaneous speed is $v_1$.
\end{itemize}
If you compute with floating-point numbers, beware of potential precision
problems.

\clearpage
\section*{Problem 3: Coloring (\begin{CJK}{UTF8}{gbsn}染色\end{CJK})
  \hfill{\normalsize\href{https://qoj.ac/problem/10196}{QOJ 10196}}}

\noindent
Problem code name: \texttt{color}.
Time limit: 1 s.
Memory limit: 512 MB.
Total points: 100.

\subsection*{Problem description}

You are given an array $A$ of $n$ positive integers, arranged in a row from
left to right.

You need to color every number in $A$ either red or blue, and then the
final score is computed as follows.

Let $C$ be an integer array of length $n$.
For every number $A_i$ in $A$ ($1 \le i \le n$):
\begin{itemize}
\item If there is no number of the same color as $A_i$ to its left, let
$C_i = 0$.
\item Otherwise, let $A_j$ be the \emph{nearest} number of the same color
to the left of $A_i$.
If $A_i = A_j$, let $C_i = A_i$; otherwise let $C_i = 0$.
\end{itemize}

Your final score is the sum of all integers in $C$, that is,
$\sum_{i=1}^{n} C_i$.
You need to maximize the final score.
Find the maximum possible final score.

\subsection*{Input}

Read the input from the file \texttt{color.in}.

\emph{This problem has multiple test cases.}

The first line of the input contains a positive integer $T$, the number of
test cases.

Then $T$ test cases follow, each in the following format.

The first line contains a positive integer $n$, the length of the array.

The second line contains $n$ positive integers $A_1, A_2, \ldots, A_n$,
the elements of the array $A$.

\subsection*{Output}

Write the output to the file \texttt{color.out}.

For each test case, output one line containing a non-negative integer, the
maximum possible value of the final score.

\subsection*{Sample 1 input}

\begin{verbatim}
3
3
1 2 1
4
1 2 3 4
8
3 5 2 5 1 2 1 4
\end{verbatim}

\subsection*{Sample 1 output}

\begin{verbatim}
1
0
8
\end{verbatim}

\subsection*{Explanation of sample 1}

For the first test case, here are three possible colorings.
\begin{enumerate}
\item Color $A_1, A_2$ red and $A_3$ blue (in the original statement the
sequence $121$ is printed with these colors).
The score is computed as follows.
\begin{itemize}
\item For $A_1$, there is no red number to its left, so $C_1 = 0$.
\item For $A_2$, the nearest red number to its left is $A_1$.
Since $A_1 \neq A_2$, $C_2 = 0$.
\item For $A_3$, there is no blue number to its left, so $C_3 = 0$.
\end{itemize}
The final score of this coloring is $C_1 + C_2 + C_3 = 0$.
\item Color all of $A_1, A_2, A_3$ red (the sequence $121$ printed all
red).
The score is computed as follows.
\begin{itemize}
\item For $A_1$, there is no red number to its left, so $C_1 = 0$.
\item For $A_2$, the nearest red number to its left is $A_1$.
Since $A_1 \neq A_2$, $C_2 = 0$.
\item For $A_3$, the nearest red number to its left is $A_2$.
Since $A_2 \neq A_3$, $C_3 = 0$.
\end{itemize}
The final score of this coloring is $C_1 + C_2 + C_3 = 0$.
\item Color $A_1, A_3$ red and $A_2$ blue (the sequence $121$ printed with
these colors).
The score is computed as follows.
\begin{itemize}
\item For $A_1$, there is no red number to its left, so $C_1 = 0$.
\item For $A_2$, there is no blue number to its left, so $C_2 = 0$.
\item For $A_3$, the nearest red number to its left is $A_1$.
Since $A_1 = A_3$, $C_3 = A_3 = 1$.
\end{itemize}
The final score of this coloring is $C_1 + C_2 + C_3 = 1$.
\end{enumerate}
It can be shown that no coloring achieves a final score greater than $1$.

For the second test case, it can be shown that every coloring has final
score $0$.

For the third test case, one optimal coloring is to color
$A_1, A_2, A_4, A_5, A_7$ red and $A_3, A_6, A_8$ blue (in the original
statement the sequence 35152124 is printed with these colors).
The corresponding array is $C = [0, 0, 0, 5, 0, 1, 2, 0]$, and the final
score is $8$.

\subsection*{Sample 2}

See the files \texttt{color/color2.in} and \texttt{color/color2.ans} in
the contestant's directory.
On QOJ these are the downloadable attachment files \texttt{ex\_color2.in}
and \texttt{ex\_color2.ans}.

\subsection*{Constraints}

For all test data it is guaranteed that $1 \le T \le 10$,
$2 \le n \le 2 \times 10^5$ and $1 \le A_i \le 10^6$.

\begin{center}
\begin{tabular}{|c|c|c|}
\hline
Test cases & $n$ & $A_i$ \\
\hline
$1 \sim 4$ & $\le 15$ & $\le 15$ \\
\hline
$5 \sim 7$ & $\le 10^2$ & $\le 10^2$ \\
\hline
$8 \sim 10$ & $\le 2{,}000$ & $\le 2{,}000$ \\
\hline
$11, 12$ & $\le 2 \times 10^4$ & $\le 10^6$ \\
\hline
$13 \sim 15$ & $\le 2 \times 10^5$ & $\le 10$ \\
\cline{1-1}\cline{3-3}
$16 \sim 20$ & & $\le 10^6$ \\
\hline
\end{tabular}
\end{center}

\clearpage
\section*{Problem 4: Arena Game
(\begin{CJK}{UTF8}{gbsn}擂台游戏\end{CJK})
  \hfill{\normalsize\href{https://qoj.ac/problem/10197}{QOJ 10197}}}

\noindent
Problem code name: \texttt{arena}.
Time limit: 1 s.
Memory limit: 512 MB.
Total points: 100.

\subsection*{Problem description}

Little S wants to hold an arena game.
If $2^k$ players take part in total, the game consists of $k$ rounds:
\begin{itemize}
\item In the first round, the players numbered $1, 2$ play one match, the
players numbered $3, 4$ play one match, and so on, up to the players
numbered $2^k - 1, 2^k$ who play one match.
\item In the second round, only the winners of the first round are kept,
and adjacent pairs play one match each, in order.
\item And so on.
In round $k - 1$, only the $4$ winners of round $k - 2$ are kept, and the
first two and the last two play one match each, which is the so-called
semifinal.
\item Round $k$ is the final between the two winners of the semifinals.
\end{itemize}

Having fixed the advancement rules, little S sets the rules of the matches
to those of a ``king of the hill'' arena.
Specifically, every player has an ability value $a_1, a_2, \cdots, a_{2^k}$,
an integer in $[0, 2^{31} - 1]$.
For every match, a number $0/1$ is first drawn by lot.
We denote the number drawn for the $G$-th match of round $R$ by $d_{R,G}$.
Drawing $0$ means that the player with the smaller number is the
defender.
Drawing $1$ means that the player with the larger number is the defender.
The defender wins if and only if his ability value satisfies $a \ge R$.
In other words, the outcome of a match depends \emph{only on the ability
value of the defender compared with the index of the current round}, and
\emph{not on the ability value of the other player}.

Now little S has received, one after another, the registrations of $n$
players, each of whom has told little S their ability value.
Little S numbers the players $1, 2, \cdots, n$ in order of registration.
Little S wants to know: after adding \emph{as few} players \emph{as
possible} so that the total number of players is a power of $2$, and after
all the players play one complete arena game, what is the \emph{sum of the
numbers} of all players who can possibly become the overall champion?

Formally, let $k$ be the smallest non-negative integer with $2^k \ge n$.
Then $(2^k - n)$ players should be added, and the ability values of the
added players can be any integers in $[0, 2^{31} - 1]$.
\emph{If an added player can possibly win, he should also be counted in the
answer.}

Of course little S finds this problem too easy, so he gives you $m$
queries $c_1, c_2, \cdots, c_m$.
For every $c_i$, little S wants you to find the answer to this problem when
only the registrations of the first $c_i$ players have been received.

\subsection*{Input}

Read the input from the file \texttt{arena.in}.

\emph{Each test file of this problem contains multiple test cases}, but the
different test cases are obtained only by modifying
$a_1, a_2, \cdots, a_n$, while everything else stays the same.
Please refer to the following format.
Here $\oplus$ denotes the bitwise XOR operator, and $a \bmod b$ denotes the
remainder of $a$ divided by $b$.

The first line of the input contains two positive integers $n, m$, the
number of registered players and the number of queries.

The second line of the input contains $n$ non-negative integers
$a'_1, a'_2, \cdots, a'_n$.
This sequence will be used to compute the real ability values.

The third line of the input contains $m$ positive integers
$c_1, c_2, \cdots, c_m$, the queries.

Let $K$ be the smallest non-negative integer with $2^K \ge n$.
In the following $K$ lines, the $R$-th line contains $2^{K-R}$ digits
(without spaces), the $G$-th of which is the number
$d_{R,G} = 0/1$ drawn by lot for the $G$-th match of round $R$.

Note that since a query only fills the number of players up to
$2^k \ge c_i$, where $k \le K$, you will not necessarily use all of the
input values.

The next line contains a positive integer $T$, the number of test cases.

Then $T$ lines follow, each describing one test case by $4$ non-negative
integers $X_0, X_1, X_2, X_3$.
The ability values of this test case are
$a_i = a'_i \oplus X_{i \bmod 4}$, where $1 \le i \le n$.

\subsection*{Output}

Write the output to the file \texttt{arena.out}.

Output $T$ lines in total.
For each test case, let $A_i$ be the answer to the $i$-th query
($1 \le i \le m$).
You only need to output one line containing one integer, the value of
$(1 \times A_1) \oplus (2 \times A_2) \oplus \cdots \oplus (m \times A_m)$.

\subsection*{Sample 1 input}

\begin{verbatim}
5 5
0 0 0 0 0
5 4 1 2 3
1001
10
1
4
2 1 0 0
1 2 1 0
0 2 3 1
2 2 0 1
\end{verbatim}

\subsection*{Sample 1 output}

\begin{verbatim}
5
19
7
1
\end{verbatim}

\subsection*{Explanation of sample 1}

There are $T = 4$ test cases; only the first one is explained here.
The real ability values of the $5$ players are $[1, 0, 0, 2, 1]$.
The $5$ queries are about the prefixes of lengths $5, 4, 1, 2, 3$,
respectively.
\begin{enumerate}
\item For the prefix of length $1$, player 1 is the only player, so the
answer is $1$.
\item For the prefix of length $2$, since $2$ is already a power of $2$, no
players need to be added.
By the lot $d_{1,1} = 1$, player 2 is the defender.
Since $a_2 < 1$, player 1 wins, and the answer is $1$.
\item For the prefix of length $3$, first, player 1 wins the match between
players 1 and 2 (because $d_{1,1} = 1$, so player 2 is the defender, and
$a_2 < 1$).
Then, although the ability value of player 4 is not yet known, player 4
surely wins the match between players 3 and 4 (because $d_{1,2} = 0$, so
player 3 is the defender, and $a_3 < 1$).
In the final between players 1 and 4, either may win (because
$d_{2,1} = 1$, so player 4 is the defender; if $a_4 < 2$ then player 1
wins, and if $a_4 \ge 2$ then player 4 wins).
Altogether, the answer is $1 + 4 = 5$.
\item For the prefix of length $4$, by the analysis of the previous item,
since $a_4 \ge 2$, the winner of the final is player 4.
\item For the prefix of length $5$, it can be shown that the players who
can possibly win are players 4, 7 and 8, so the answer is $19$.
\end{enumerate}
Therefore the answer for this test case is
$(1 \times 19) \oplus (2 \times 4) \oplus (3 \times 1) \oplus (4 \times 1)
\oplus (5 \times 5) = 5$.

\subsection*{Sample 2}

See the files \texttt{arena/arena2.in} and \texttt{arena/arena2.ans} in
the contestant's directory.
On QOJ these are the downloadable attachment files \texttt{ex\_arena2.in}
and \texttt{ex\_arena2.ans}.

This sample satisfies special property A.

\subsection*{Sample 3}

See the files \texttt{arena/arena3.in} and \texttt{arena/arena3.ans} in
the contestant's directory.
On QOJ these are the downloadable attachment files \texttt{ex\_arena3.in}
and \texttt{ex\_arena3.ans}.

This sample satisfies special property B.

\subsection*{Sample 4}

See the files \texttt{arena/arena4.in} and \texttt{arena/arena4.ans} in
the contestant's directory.
On QOJ these are the downloadable attachment files \texttt{ex\_arena4.in}
and \texttt{ex\_arena4.ans}.

\subsection*{Sample 5}

See the files \texttt{arena/arena5.in} and \texttt{arena/arena5.ans} in
the contestant's directory.
On QOJ these are the downloadable attachment files \texttt{ex\_arena5.in}
and \texttt{ex\_arena5.ans}.

\subsection*{Constraints}

For all test data it is guaranteed that $2 \le n, m \le 10^5$,
$0 \le a_i, X_j < 2^{31}$, $1 \le c_i \le n$ and $1 \le T \le 256$.

\begin{center}
\small
\begin{tabular}{|c|c|c|c|c|}
\hline
Test cases & $T =$ & $n, m \le$ & Special property A & Special property B \\
\hline
$1 \sim 3$ & $1$ & $8$ & no & no \\
\cline{1-1}\cline{3-4}
$4, 5$ & & $500$ & yes & \\
\cline{1-1}\cline{4-5}
$6 \sim 8$ & & & no & yes \\
\cline{1-1}\cline{3-3}\cline{5-5}
$9, 10$ & & $5{,}000$ & & no \\
\cline{1-1}\cline{3-4}
$11, 12$ & & $10^5$ & yes & \\
\cline{1-1}\cline{4-5}
$13 \sim 15$ & & & no & yes \\
\cline{1-2}\cline{5-5}
$16, 17$ & $4$ & & & no \\
\cline{1-2}
$18, 19$ & $16$ & & & \\
\cline{1-2}
$20, 21$ & $64$ & & & \\
\cline{1-2}
$22, 23$ & $128$ & & & \\
\cline{1-2}
$24, 25$ & $256$ & & & \\
\hline
\end{tabular}
\end{center}

Special property A: it is guaranteed that every queried $c_i$ is a power of
$2$.

Special property B: it is guaranteed that all $d_{R,G} = 0$.

\section*{Source}

The statements were taken from the QOJ pages of the contest CSP-S 2024
(the statement PDF embedded in each page), in contest order:
\begin{itemize}
\item Problem 1 (duel): \texttt{https://qoj.ac/problem/10194}
\item Problem 2 (detect): \texttt{https://qoj.ac/problem/10195}
\item Problem 3 (color): \texttt{https://qoj.ac/problem/10196}
\item Problem 4 (arena): \texttt{https://qoj.ac/problem/10197}
\end{itemize}

\end{document}
